\section{Direct and Semidirect Products and Abelian Groups}

\subsection{Direct Products}

\begin{exercise}{5.1.1}
    Show that the center of a direct product is the direct product of the centers:
    \[
    	Z(G_1 \times G_2 \times \cdots \times G_n) = Z(G_1) \times Z(G_2) \times \cdots \times Z(G_n)
    \]
    Deduce that the direct product of two groups is Abelian if and only if each of the factors is Abelian.
\end{exercise}

\begin{sol}
    Let $G = G_1 \times G_2 \times \cdots \times G_n$.
	Let $z = (z_1, z_2, \ldots, z_n) \in Z(G)$.
	Then for any $g = (g_1, g_2, \ldots, g_n) \in G$, we have $zg = gz$, or equivalently,
    \[
    	(z_1g_1, z_2g_2, \ldots, z_ng_n) = (g_1z_1, g_2z_2, \ldots, g_nz_n)
    \]
    This implies that $z_i g_i = g_i z_i$ for all $i$, so $z_i \in Z(G_i)$ for all $i$.
	Thus, $z \in Z(G_1) \times Z(G_2) \times \cdots \times Z(G_n)$, and we have shown that $Z(G) \subseteq Z(G_1) \times Z(G_2) \times \cdots \times Z(G_n)$.

    Conversely, let $z = (z_1, z_2, \ldots, z_n) \in Z(G_1) \times Z(G_2) \times \cdots \times Z(G_n)$.
	Then for any $g = (g_1, g_2, \ldots, g_n) \in G$, we have
    \[
    	zg = (z_1g_1, z_2g_2, \ldots, z_ng_n) = (g_1z_1, g_2z_2, \ldots, g_nz_n) = gz
    \]
    Thus, $z \in Z(G)$, and we have shown that $Z(G_1) \times Z(G_2) \times \cdots \times Z(G_n) \subseteq Z(G)$.

    In particular, if two groups $G_1$ and $G_2$ are Abelian, then $Z(G_1) = G_1$ and $Z(G_2) = G_2$, so $Z(G_1 \times G_2) = Z(G_1) \times Z(G_2) = G_1 \times G_2$, which is Abelian.
    Conversely, if $G_1 \times G_2$ is Abelian, then $Z(G_1 \times G_2) = G_1 \times G_2$, so $Z(G_1) \times Z(G_2) = G_1 \times G_2$, which implies that $Z(G_1) = G_1$ and $Z(G_2) = G_2$, so both $G_1$ and $G_2$ are Abelian.
\end{sol}

\begin{exercise}{5.1.2}
    Let $G_1, G_2, \ldots, G_n$ be groups and $G = G_1 \times \cdots \times G_n$.
	Let $I$ be a proper, nonempty subset of $\set{1, \ldots, n}$ and let $J = \set{1, \ldots, n} \setminus I$.
	Define $G_I$ to be the set of elements of $G$ that have the identity $G_j$ in position $j$ for all $j \in J$.
    \begin{subexercises}
        \item Prove that $G_I$ is isomorphic to the direct product of the $G_i$ for $i \in I$.
        \item Prove that $G_I$ is a normal subgroup of $G$ and that $G/G_I \cong G_J$.
        \item Prove that $G \cong G_I \times G_J$.
    \end{subexercises}
\end{exercise}

\begin{solenum}
    \item Define the map
    \[
        \phi : G_I \to \prod_{i \in I} G_i, \quad \phi((g_1, g_2, \ldots, g_n)) = (g_i)_{i \in I}.
    \]
    We have $\phi$ is a homomorphism since the group operation in $G_I$ is componentwise, and the group operation in $\prod_{i \in I} G_i$ is also componentwise.
    To see that $\phi$ is injective, suppose $\phi((g_1, \ldots, g_n)) = (1)_{i \in I}$, the identity in $\prod_{i \in I} G_i$.
    Then $g_i = 1$ for all $i \in I$, and since $g_j = 1$ for all $j \in J$ by definition of $G_I$, we have $(g_1, \ldots, g_n) = (1, \ldots, 1)$, the identity in $G_I$.

    To see that $\phi$ is surjective, let $(h_i)_{i \in I} \in \prod_{i \in I} G_i$.
    Then $(g_1, \ldots, g_n) \in G_I$ defined by $g_i = h_i$ for $i \in I$ and $g_j = 1$ for $j \in J$ satisfies $\phi((g_1, \ldots, g_n)) = (h_i)_{i \in I}$.
    Hence, $\phi$ is surjective, and we conclude that $\phi$ is an isomorphism.
    \item As shown in part (a), we may identify $G_J$ with the direct product of the $G_j$ for $j \in J$.
    Define the map
    \[
        \pi_J : G \to G_J, \quad \pi_J((g_1, g_2, \ldots, g_n)) = (g_j)_{j \in J}.
    \]
    As in the previous part, $\pi_J$ is seen to be a surjective homomorphism.
    The kernel of $\pi_J$ consists of all elements $(g_1, \ldots, g_n) \in G$ such that $g_j = 1$ for all $j \in J$, which is precisely the definition of $G_I$, so $\ker \pi_J = G_I$.
    By the First Isomorphism Theorem, we have $G/G_I \cong \pi_J(G) = G_J$.
    \item Using part (a) to identify $G_I$ and $G_J$ with their respective direct products, define the map
    \[
        \psi : G \to G_I \times G_J, \quad \psi((g_1, g_2, \ldots, g_n)) = ((g_i)_{i \in I}, (g_j)_{j \in J}).
    \]
    It is clear that $\psi$ is a homomorphism, and it is injective since if $\psi((g_1, \ldots, g_n)) = ((1)_{i \in I}, (1)_{j \in J})$, then $g_i = 1$ for all $i \in I$ and $g_j = 1$ for all $j \in J$, so $(g_1, \ldots, g_n) = (1, \ldots, 1)$.

    It is surjective since for any $((h_i)_{i \in I}, (k_j)_{j \in J}) \in G_I \times G_J$, we have $(g_1, \ldots, g_n) \in G$ defined by $g_i = h_i$ for $i \in I$ and $g_j = k_j$ for $j \in J$ satisfies $\psi((g_1, \ldots, g_n)) = ((h_i)_{i \in I}, (k_j)_{j \in J})$.
    Hence, $\psi$ is an isomorphism, and we conclude that $G \cong G_I \times G_J$.
\end{solenum}

\begin{exercise}{5.1.3}
    Under the notation of the preceding exercise, let $I$ and $K$ be any disjoint nonempty subsets of $\set{1, 2, \ldots, n}$, and let $G_I$ and $G_K$ be the subgroups of $G$ defined above.
	Prove that $xy = yx$ for all $x \in G_I$ and $y \in G_K$.
\end{exercise}

\begin{sol}
    Let $x \in G_I$ and $y \in G_K$.
	Then $x$ has the identity in all coordinates corresponding to $J = \set{1, 2, \ldots, n} - I$, and $y$ has the identity in all coordinates corresponding to $L = \set{1, 2, \ldots, n} - K$.
	Since $I$ and $K$ are disjoint, we have $J \cup L = \set{1, 2, \ldots, n}$, so every coordinate of $x$ and $y$ is the identity in at least one of them.
	Thus, when we compute the product $xy$, we have
    \[
    	xy = (x_1y_1, x_2y_2, \ldots, x_ny_n) = (y_1x_1, y_2x_2, \ldots, y_nx_n) = yx
    \]
    as desired.
\end{sol}

\begin{exercise}{5.1.4}
    Let $A$ and $B$ be finite groups, and let $p$ be a prime.
	Prove that any Sylow $p$-subgroup of $A \times B$ is of the form $P \times Q$, where $P \in \syl_p(A)$ and $Q \in \syl_p(B)$.
	Prove that $n_p(A \times B) = n_p(A)n_p(B)$.
	Generalize both of these results to a direct product of any finite number of finite groups (so that the number of Sylow $p$-subgroups of a direct product is the product of the numbers of Sylow $p$-subgroups of the factors).
\end{exercise}

\begin{sol}
    Let $|A| = p^\alpha m$ and $|B| = p^\beta n$, where $p \nmid m$ and $p \nmid n$.
	Then $|A \times B| = |A||B| = p^{\alpha + \beta}mn$.
	Let $R \in \syl_p(A \times B)$.
	Then $|R| = p^{\alpha + \beta}$, so $R$ is a Sylow $p$-subgroup of $A \times B$.
	Let $\pi_A : A \times B \to A$ and $\pi_B : A \times B \to B$ be the projection maps.
	Then $\pi_A(R)$ is a $p$-subgroup of $A$, so $\pi_A(R) \subseteq P$ for some $P \in \syl_p(A)$.
	Similarly, $\pi_B(R) \subseteq Q$ for some $Q \in \syl_p(B)$.
	Thus, $R \subseteq P \times Q$.
	Since $|P \times Q| = |P||Q| = p^{\alpha + \beta}$, we must have $R = P \times Q$.

    To show $n_p(A \times B) = n_p(A)n_p(B)$, we note that the Sylow $p$-subgroups of $A \times B$ are precisely the subgroups of the form $P \times Q$, where $P \in \syl_p(A)$ and $Q \in \syl_p(B)$.
	Since there are $n_p(A)$ choices for $P$ and $n_p(B)$ choices for $Q$, we have $n_p(A \times B) = n_p(A)n_p(B)$.

	Generalization to a direct product of any finite number of finite groups follows by induction on the number of factors.
\end{sol}

\begin{exercise}{5.1.5}
    Exhibit a non-normal subgroup of $Q_8 \times Z_4$ (note that every subgroup of each factor is normal).
\end{exercise}

\begin{sol}
    Note that $Q_8$ is non-Abelian, so we may use the idea of a diagonal subgroup from \exref{3.1.38}.
	In particular, let $H = \gen{(i, 1)} = \set{(1, 0), (i, 1), (-1, 2), (-i, 3)}$.
	Then $H$ is a subgroup of $Q_8 \times Z_4$ that is isomorphic to $\gen{i}$.
	However,
    \[
    	(j, 0)(i, 1)(j, 0)\inv = (ji(-j), 0 + 1 + 0) = (-i, 1)
    \]
    is not in $H$, so $H$ is not normal in $Q_8 \times Z_4$.
\end{sol}

\begin{exercise}{5.1.6}
    Show that all subgroups of $Q_8 \times E_{2^n}$ are normal.
\end{exercise}

\begin{sol}
    Let $H \leq Q_8 \times E_{2^n}$.
	Suppose $(a, b) \in H$ and $(x, y) \in Q_8 \times E_{2^n}$.
	Note that $(a\inv, b) \in H$, since every $b \in E_{2^n}$ is its own inverse.
	Then
    \[
    	(x, y)(a, b)(x, y)\inv = (xax\inv, yby\inv) = (xax\inv, b),
    \]
    where we used the fact that $E_{2^n}$ is Abelian.
	Since every subgroup of $Q_8$ is normal, we have $xax\inv \in \gen a$.
    Note that $xax\inv \neq a^2$, since $a^2$ is the unique element of order 2 in $Q_8$, and $xax\inv$ has the same order as $a$.
    Then either $xax\inv = a$ or $xax\inv = a\inv$.
    Both $(a, b)$ and $(a\inv, b)$ are in $H$, so $(xax\inv, b) \in H$.
    Hence, $H \nsub Q_8 \times E_{2^n}$.
\end{sol}

\begin{exercise}{5.1.7}
    Let $G_1, G_2, \ldots, G_n$ be groups and $\pi$ be a fixed element of $S_n$.
	Prove that the map 
    \[
    	\phi_\pi : G_1 \times G_2 \times \cdots \times G_n \to G_{\pi\inv(1)} \times G_{\pi\inv(2)} \times \cdots \times G_{\pi\inv(n)}
    \]
    defined by
    \[
    	\phi_\pi(g_1, g_2, \ldots, g_n) = (g_{\pi\inv(1)}, g_{\pi\inv(2)}, \ldots, g_{\pi\inv(n)})
    \]
    is an isomorphism (so that changing the order of the factors in a direct product does not change the isomorphism type).
\end{exercise}

\begin{sol}
    Let $\pi \in S_n$ and let $\phi_\pi$ be the map defined in the problem statement.
	Let $g = (g_1, g_2, \ldots, g_n)$ and $h = (h_1, h_2, \ldots, h_n)$ be elements of $G_1 \times G_2 \times \cdots \times G_n$.
	Then
    \begin{align*}
        \phi_\pi(gh) & = \phi_\pi(g_1h_1, g_2h_2, \ldots, g_nh_n) \\
        & = (g_{\pi\inv(1)}h_{\pi\inv(1)}, g_{\pi\inv(2)}h_{\pi\inv(2)}, \ldots, g_{\pi\inv(n)}h_{\pi\inv(n)}) \\
        & = (g_{\pi\inv(1)}, g_{\pi\inv(2)}, \ldots, g_{\pi\inv(n)})(h_{\pi\inv(1)}, h_{\pi\inv(2)}, \ldots, h_{\pi\inv(n)}) \\
        & = \phi_\pi(g)\phi_\pi(h)
    \end{align*}
    Noting that $\pi \in S_n$ is a bijection, consider the map $\phi_{\pi\inv} : G_{\pi\inv(1)} \times G_{\pi\inv(2)} \times \cdots \times G_{\pi\inv(n)} \to G_1 \times G_2 \times \cdots \times G_n$ defined by
    \[
    	\phi_{\pi\inv}(k_1, k_2, \ldots, k_n) = (k_{\pi(1)}, k_{\pi(2)}, \ldots, k_{\pi(n)})
    \]
    Then $\phi_{\pi\inv}$ is a homomorphism by the same argument as above.
	Moreover, $\phi_{\pi\inv}$ is the inverse of $\phi_\pi$, since
    \[
    	\phi_\pi(\phi_{\pi\inv}(k_1, k_2, \ldots, k_n)) = \phi_\pi(k_{\pi(1)}, k_{\pi(2)}, \ldots, k_{\pi(n)}) = (k_1, k_2, \ldots, k_n)
    \]
    and
    \[
    	\phi_{\pi\inv}(\phi_\pi(g_1, g_2, \ldots, g_n)) = \phi_{\pi\inv}(g_{\pi\inv(1)}, g_{\pi\inv(2)}, \ldots, g_{\pi\inv(n)}) = (g_1, g_2, \ldots, g_n).
    \]
    Therefore, $\phi_\pi$ is an isomorphism.
\end{sol}

\begin{exercise}{5.1.8}
    Let $G_1 = G_2 = \cdots = G_n$ and let $G = G_1 \times \cdots \times G_n$.
	Under the notation of the preceding exercise show that $\phi_\pi \in \text{Aut}(G)$.
	Show also that the map $\pi \mapsto \phi_\pi$ is an injective homomorphism of $S_n$ into $\text{Aut}(G)$. (In particular, $\phi_{\pi_1} \circ \phi_{\pi_2} = \phi_{\pi_1 \pi_2}$.
	It is at this point that the $\pi^{-1}$'s in the definition of $\phi_\pi$ are needed.
	The underlying reason for this is that if $e_i$ is the $n$-tuple with 1 in position $i$ and zeros elsewhere, $1 \leq i \leq n$, then $S_n$ acts on $\set{e_1, \ldots, e_n}$ by $\pi \cdot e_i = e_{\pi(i)}$; this is a left group action.
	If the $n$-tuple $(g_1, \ldots, g_n)$ is represented by $g_1 e_1 + \cdots + g_n e_n$, then this left group action on $\set{e_1, \ldots, e_n}$ extends to a left group action on sums by
    \[
    	\pi \cdot (g_1 e_1 + g_2 e_2 + \cdots + g_n e_n) = g_1 e_{\pi(1)} + g_2 e_{\pi(2)} + \cdots + g_n e_{\pi(n)}.
    \]
    The coefficient of $e_{\pi(i)}$ on the right-hand side is $g_i$, so the coefficient of $e_i$ is $g_{\pi^{-1}(i)}$.
	Thus, the right-hand side may be rewritten as $g_{\pi^{-1}(1)} e_1 + g_{\pi^{-1}(2)} e_2 + \cdots + g_{\pi^{-1}(n)} e_n$, which is precisely the sum attached to the $n$-tuple $(g_{\pi^{-1}(1)}, g_{\pi^{-1}(2)}, \ldots, g_{\pi^{-1}(n)})$.
	In other words, any permutation of the ``position vectors'' $e_1, \ldots, e_n$ (which fixes their coefficients) is the same as the inverse permutation on the coefficients (fixing the $e_i$'s).
	If one uses $\pi$'s in place of $\pi^{-1}$'s in the definition of $\phi_\pi$ then the map $\pi \mapsto \phi_\pi$ is not necessarily a homomorphism--it corresponds to a \emph{right} group action.)
\end{exercise}

\begin{sol}
    Note that $\phi_\pi \in \aut(G)$ since $\phi_\pi$ is an isomorphism from $G$ to itself.
	Let $\pi_1, \pi_2 \in S_n$.
	Then
    \begin{align*}
        \phi_{\pi_1}(\phi_{\pi_2}(g_1, g_2, \ldots, g_n)) & = \phi_{\pi_1}(g_{\pi_2\inv(1)}, g_{\pi_2\inv(2)}, \ldots, g_{\pi_2\inv(n)}) \\
        & = (g_{\pi_2\inv(\pi_1\inv(1))}, g_{\pi_2\inv(\pi_1\inv(2))}, \ldots, g_{\pi_2\inv(\pi_1\inv(n))}) \\
        & = (g_{(\pi_1 \pi_2)\inv(1)}, g_{(\pi_1 \pi_2)\inv(2)}, \ldots, g_{(\pi_1 \pi_2)\inv(n)}) \\
        & = \phi_{\pi_1 \pi_2}(g_1, g_2, \ldots, g_n)
    \end{align*}
    Therefore, the map $\pi \mapsto \phi_\pi$ is a homomorphism.
	To see that it is injective, note that if $\phi_\pi$ is the identity automorphism, then for all $i$, $g_{\pi\inv(i)} = g_i$ for all $g_i \in G_i$, which implies $\pi$ is the identity permutation.
	Hence, the map is injective.
\end{sol}

\begin{exercise}{5.1.9}
    Let $G_i$ be a field $F$ for all $i$ and use the preceding exercise to show that the set of $n \times n$ matrices with one 1 in each row and each column is a subgroup of $\gl_n(F)$ isomorphic to $S_n$ (these matrices are the \emph{permutation matrices} since they simply permute the standard basis $e_1, \ldots, e_n$ (as above) of the $n$-dimensional vector space $F \times F \times \cdots \times F$).
\end{exercise}

\begin{sol}
    Let $G = F \times F \times \cdots \times F$ ($n$ times) and let $\pi \in S_n$.
	Then $\phi_\pi$ is an automorphism of $G$ that permutes the standard basis vectors $e_1, e_2, \ldots, e_n$.
	In particular, $\phi_\pi(e_i) = e_{\pi(i)}$ for all $i$.
	Thus, the matrix representation of $\phi_\pi$ with respect to the standard basis is a permutation matrix.
	Since the map $\pi \mapsto \phi_\pi$ is an injective homomorphism from $S_n$ to $\aut(G)$, the set of permutation matrices forms a subgroup of $\gl_n(F)$ that is isomorphic to $S_n$.
\end{sol}

\begin{exercise}{5.1.10}
    Let $p$ be a prime.
	Let $A$ and $B$ be two cyclic groups of order $p$ with generators $x$ and $y$ respectively.
	Set $E = A \times B$ so that $E$ is the elementary Abelian group of order $p^2$: $E_{p^2}$ Prove that the distinct subgroups of $E$ of order $p$ are
    \[
    	\langle x \rangle, \quad \langle xy \rangle, \quad \langle xy^2 \rangle, \quad \ldots, \quad \langle xy^{p-1} \rangle, \quad \langle y \rangle
    \]
	(note that there are $p + 1$ of them).
\end{exercise}

\begin{sol}
    Note that any proper, non-trivial subgroup of $E$ must have order $p$, and every element of $E$ has order $p$.
	The text shows that there are $p + 1$ subgroups of $E$ of order $p$, so it suffices to show that the subgroups listed in the problem statement are distinct.
	Note that $\gen x$ and $\gen y$ are distinct since they have different generators.
	Moreover, $\gen{xy^i}$ is distinct from $\gen x$ and $\gen y$ for all $i$, since $xy^i$ is not a power of $x$ or $y$ (as can be seen because $\gen x \cap \gen y$ is trivial).
	Finally, $\gen{xy^i}$ is distinct from $\gen{xy^j}$ for $i \neq j$, since if $\gen{xy^i} = \gen{xy^j}$, $xy^i$ is a power of $xy^j$, which implies $y^{i-j}$ is a power of $x$, which is impossible since $\gen x \cap \gen y$ is trivial.
	Thus, all the subgroups listed in the problem statement are distinct.
\end{sol}

\begin{exercise}{5.1.11}
    Let $p$ be a prime, and let $n \in \zp$.
	Find a formula for the number of subgroups of order $p$ in the elementary Abelian group $E_{p^n}$.
\end{exercise}

\begin{sol}
    Every element of $E_{p^n}$ has order $p$, so every subgroup of order $p$ is generated by a single non-identity element.
	Since there are $p^n - 1$ non-identity elements in $E_{p^n}$, and each subgroup of order $p$ has $p - 1$ generators (the non-identity elements of the subgroup), the number of subgroups of order $p$ in $E_{p^n}$ is
    \[
    	\frac{p^n - 1}{p - 1} = 1 + p + p^2 + \cdots + p^{n-1}
    \]
    as desired.
\end{sol}

\begin{spex}{5.1.12}
    Let $A$ and $B$ be groups.
	Assume $Z(A)$ contains a subgroup $Z_1$ and $Z(B)$ contains a subgroup $Z_2$ with $Z_1 \cong Z_2$.
	Let this isomorphism be given by the map $x_i \mapsto y_i$ for all $x_i \in Z_1$.
	A \textbf{\emph{central product}} of $A$ and $B$ is a quotient
    \[
    	(A \times B)/Z \quad \text{where} \quad Z = \set{(x_i, y_i^{-1}) \mid x_i \in Z_1}
    \]
    and is denoted by $A \ast B$--it is not unique since it depends on $Z_1$, $Z_2$ and the isomorphism between them. (Think of $A \ast B$ as the direct product of $A$ and $B$ ``collapsed'' by identifying each element $x_i \in Z_1$ with its corresponding element $y_i \in Z_2$.)
    \begin{subexercises}
        \item Prove that the images of $A$ and $B$ in the quotient group $A \ast B$ are isomorphic to $A$ and $B$, respectively, and that these images intersect in a central subgroup isomorphic to $Z_1$.
	    Find $|A \ast B|$.
        \item Let $Z_4 = \langle x \rangle$.
	    Let $D_8 = \langle r, s \rangle$ and $Q_8 = \langle i, j \rangle$ be given by their usual generators and relations.
	    Let $Z_4 \ast D_8$ be the central product of $Z_4$ and $D_8$ which identifies $x^2$ and $r^2$ (i.e., $Z_1 = \langle x^2 \rangle$, $Z_2 = \langle r^2 \rangle$ and the isomorphism is $x^2 \mapsto r^2$) and let $Z_4 \ast Q_8$ be the central product of $Z_4$ and $Q_8$ which identifies $x^2$ and $-1$.
	    Prove that $Z_4 \ast D_8 \cong Z_4 \ast Q_8$.
    \end{subexercises}
\end{spex}

\begin{solenum}
    \item Consider the maps $\pi_A : A \to A \ast B$ and $\pi_B : B \to A \ast B$ defined by $\pi_A(a) = (a, 1)Z$ and $\pi_B(b) = (1, b)Z$.
	It is clear that $\pi_A$ and $\pi_B$ are homomorphisms.
	Moreover, $\ker\pi_A$ consists of the elements $a \in A$ such that $(a, 1) \in Z$.
	Then $y_i = 1$, and since $x_i \mapsto y_i$ is an isomorphism, we have $x_i = 1$, so $a = 1$.
	Thus, $\ker\pi_A$ is trivial, and $\pi_A$ is injective.
	Similarly, $\ker\pi_B$ is trivial, and $\pi_B$ is injective.
	Therefore, the images of $A$ and $B$ in $A \ast B$ are isomorphic to $A$ and $B$, respectively.
	Moreover, the images of $A$ and $B$ in $A \ast B$ intersect in the subgroup $\set{(x_i, 1)Z \mid x_i \in Z_1}$, which is isomorphic to $Z_1$.
	Finally, since $Z$ has $|Z_1|$ elements, we have
    \[
    	|A \ast B| = \frac{|A||B|}{|Z_1|}.
    \]
    \item Denote $Z_D$ and $Z_Q$ as the subgroups of $Z_4 \times D_8$ and $Z_4 \times Q_8$ that are identified in the central products.
    Put $x$ and $y$ as the generators of the $Z_4$ factors in $Z_4 \ast D_8$ and $Z_4 \ast Q_8$, respectively.
    Put $r$ and $s$ as the generators of $D_8$ and $i$ and $j$ as the generators of $Q_8$.
    Let $\b x = (x, 1)Z_D$ be the coset of $Z_D$ in $Z_4 \ast D_8$, and denote all other elements of $Z_4 \ast D_8$ and $Z_4 \ast Q_8$ similarly.

    A preliminary observation is that we need an element in $Z_4 \ast Q_8$ that corresponds to $\b s \in Z_4 \ast D_8$.
    We claim that the element $\b z = \b j \b y\inv$ has order 2 and conjugates $\b i$ to $\b i\inv$.
    Note that $|\b z| \neq 1$, for otherwise $\b j = \b y$. This would imply $\b j \in \pi_{Z_4}(y) = \gen{\b y}$, but $\b j$ does not commute with $\b i$, so $\b j \notin \gen{\b y}$. Then
    \[
        \b z^2 = \b j \b y\inv \b j \b y\inv = \b j^2 \b y^{-2} = \b 1,
    \]
    so $\b z$ has order 2.
    Moreover, we have
    \[
        \b z \b i \b z\inv = \b j \b y\inv \b i \b y \b j = \b j \b i \b j = \b i\inv.
    \]
    Then $\b z$ has the same properties as $\b s$, so we may define a map $\phi : Z_4 \ast D_8 \to Z_4 \ast Q_8$ by
    \[
        \phi(\b x^\alpha \b r^\beta \b s^\gamma) = \b y^\alpha \b i^\beta \b z^\gamma.
    \]

    We claim that $\phi$ is well-defined.
    To see this, note that the relations in $Z_4 \ast D_8$ are satisfied in $Z_4 \ast Q_8$ under the map $\phi$.
    In particular, we have in $Z_4 \ast D_8$ that $\b r^4 = \b s^2 = \b 1$, $\b x^2 = \b r^2$, and $\b s\b r\b s = \b r\inv$.
    Then in $Z_4 \ast Q_8$, we have
    \[
        \phi(\b r^4) = \b i^4 = \b 1, \quad \phi(\b s^2) = \b z^2 = \b 1, \quad \phi(\b x^2) = \b y^2 = \b i^2 = \phi(\b r^2), \quad \text{and} \quad \phi(\b s\b r\b s) = \b z\b i\b z = \b i\inv.
    \]
    Additionally, $\b y$ is central in $Z_4 \ast Q_8$, which follows from the fact that $Z_4$ is central in $Z_4 \ast D_8$ and $\phi$ preserves the relations.
    Thus, $\phi$ is a homomorphism.

    To see that $\phi$ is surjective, note that any element of $Z_4 \ast Q_8$ can be written as $\b y^\alpha \b i^\beta \b z^\gamma$, which is the image of $\b x^\alpha \b r^\beta \b s^\gamma$ under $\phi$.
    Since $Z_4$ and $Q_8$ are generated by $\b y$ and $\b i, \b j$ respectively, putting $\b j = \b z \b y$ shows that $\b y, \b i, \b z$ generate $Z_4 \ast Q_8$.
    From part (a), we know that $|Z_4 \ast D_8| = |Z_4||D_8|/|Z_1| = 16$ and $|Z_4 \ast Q_8| = |Z_4||Q_8|/|Z_1| = 16$, so $\phi$ is a surjective homomorphism between two groups of the same order, and hence it is an isomorphism.
\end{solenum}

\begin{exercise}{5.1.13}
    Give presentations for the groups $Z_4 \ast D_8$ and $Z_4 \ast Q_8$ constructed in the preceding exercise.
\end{exercise}

\begin{sol}
    To construct the presentations, we note that $\b x$ should satisfy the relations of $Z_4$, $\b r$ and $\b s$ should satisfy the relations of $D_8$, and $\b i$ and $\b z$ should satisfy the relations of $Q_8$.
	Moreover, we have the relations $\b x^2 = \b r^2$ in $Z_4 \ast D_8$, since $\b x^2\b r^{-2} = \b 1$.
	Lastly, $\b x$ is central in that $\b x \b r = \b r\b x$ and $\b x\b s = \b s\b x$ in $Z_4 \ast D_8$.
	Then a presentation for $Z_4 \ast D_8$ is
    \[
    	\gen{\b x, \b r, \b s \mid \b r^4 = \b s^2 = \b 1, \b x^2 = \b r^2, \b s\b r\b s = \b r\inv, \b x\b r = \b r\b x, \b x\b s = \b s\b x}.
    \]
    To find a presentation for $Z_4 \ast Q_8$, we may use the isomorphism $\phi$ from the preceding exercise to rewrite the presentation for $Z_4 \ast D_8$ in terms of $\b y, \b i, \b z$.
	In particular, we have $\b y^2 = \b i^2\b z^2$, $\b z^2 = \b 1$, and $\b y\b i = \b i\b y\b z$.
	Thus, a presentation for $Z_4 \ast Q_8$ is
    \[
    	\gen{\b y, \b i, \b z \mid \b i^4 = \b z^2 = \b 1, \b y^2 = \b i^2, \b z\b i \b z = \b i\inv, \b y\b i = \b i\b y, \b y\b z = \b z\b y}. \qedhere
    \]
\end{sol}

\begin{exercise}{5.1.14}
    Let $G = A_1 \times A_2 \times \cdots \times A_n$, and for each $i$, let $B_i$ be a normal subgroup of $A_i$.
	Prove that $B_1 \times B_2 \times \cdots \times B_n \nsub G$, and that
    \[
    	(A_1 \times A_2 \times \cdots \times A_n)/(B_1 \times B_2 \times \cdots \times B_n) \cong (A_1/B_1) \times (A_2/B_2) \times \cdots \times (A_n/B_n).
    \]
\end{exercise}

\begin{sol}
    Let $G = A_1 \times A_2 \times \cdots \times A_n$ and let $B_i \nsub A_i$ for all $i$.
	Let $H = B_1 \times B_2 \times \cdots \times B_n$.
	It is clear that $H$ is a subgroup of $G$.
	Moreover, since each $B_i$ is normal in $A_i$, we have that for any $g = (a_1, a_2, \ldots, a_n) \in G$ and any $h = (b_1, b_2, \ldots, b_n) \in H$, we have
    \[
    	ghg\inv = (a_1b_1a_1\inv, a_2b_2a_2\inv, \ldots, a_nb_na_n\inv) \in H
    \]
    since $a_ib_ia_i\inv \in B_i$ for all $i$.
	Thus, $H$ is normal in $G$.

    Consider the map $\phi : G \to (A_1/B_1) \times (A_2/B_2) \times \cdots \times (A_n/B_n)$ defined by
    \[
    	\phi((a_1, a_2, \ldots, a_n)) = (a_1B_1, a_2B_2, \ldots, a_nB_n).
    \]
    It is clear that $\phi$ is a homomorphism.
	Moreover, $\phi$ is surjective since for any $(a_1B_1, a_2B_2, \ldots, a_nB_n)$, we have $\phi((a_1, a_2, \ldots, a_n)) = (a_1B_1, a_2B_2, \ldots, a_nB_n)$.
	Finally, the kernel of $\phi$ is the set of all $(a_1, a_2, \ldots, a_n)$ such that $a_i \in B_i$ for all $i$, which is exactly $H$.
	Thus, by the First Isomorphism Theorem, we have
    \[
    	G/H \cong (A_1/B_1) \times (A_2/B_2) \times \cdots \times (A_n/B_n)
    \]
    as desired.
\end{sol}

\begin{exercise}{5.1.15}
    Let $I$ be any nonempty index set and let $(G_i, \star_i)$ be a group for each $i \in I$.
	The \emph{direct product} of the groups $G_i$, $i \in I$ is the set $G = \prod_{i \in I} G_i$ (the Cartesian product of the $G_i$'s) with a binary operation defined as follows: if $\prod a_i$ and $\prod b_i$ are elements of $G$, then their product in $G$ is given by
    \[
    	\left(\prod_{i \in I} a_i\right)\left(\prod_{i \in I} b_i\right) = \prod_{i \in I} (a_i \star_i b_i)
    \]
    (i.e., the group operation in the direct product is defined componentwise).
    \begin{subexercises}
        \item Show that this binary operation is well-defined and associative.
        \item Show that the element $\prod 1_i$ satisfies the axiom for the identity of $G$, where $1_i$ is the identity of $G_i$ for all $i$.
        \item Show that the element $\prod a_i^{-1}$ is the inverse of $\prod a_i$, where the inverse of each component element $a_i$ is taken in the group $G_i$.
    \end{subexercises}
    Conclude that the direct product is a group. (Note that if $I = \set{1, 2, \ldots, n}$, this definition of the direct product is the same as the $n$-tuple definition in the text.)
\end{exercise}

\begin{solenum}
    \item To show that the operation is well-defined, let $\prod \alpha_i = \prod a_i$ and $\prod \beta_i = \prod b_i$.
	Then $\alpha_i = a_i$ and $\beta_i = b_i$ for all $i$, so
    \[
    	\left(\prod_{i \in I} \alpha_i\right)\left(\prod_{i \in I} \beta_i\right) = \prod_{i \in I} (\alpha_i \star_i \beta_i) = \prod_{i \in I} (a_i \star_i b_i)
    \]
    as desired.
	To show associativity, let $\prod a_i$, $\prod b_i$, and $\prod c_i$ be elements of $G$.
	Then
    \begin{align*}
        \left(\prod_{i \in I} a_i\right)\left(\left(\prod_{i \in I} b_i\right)\left(\prod_{i \in I} c_i\right)\right) & = \left(\prod_{i \in I} a_i\right)\left(\prod_{i \in I} (b_i \star_i c_i)\right) \\
        & = \prod_{i \in I} (a_i \star_i (b_i \star_i c_i)) \\
        & = \prod_{i \in I} ((a_i \star_i b_i) \star_i c_i) \\
        & = \left(\prod_{i \in I} (a_i \star_i b_i)\right)\left(\prod_{i \in I} c_i\right) \\
        & = \left(\left(\prod_{i \in I} a_i\right)\left(\prod_{i \in I} b_i\right)\right)\left(\prod_{i \in I} c_i\right)
    \end{align*}
    as desired.
    \item Let $\prod a_i$ be an element of $G$.
	Then
    \[
    	\left(\prod_{i \in I} 1_i\right)\left(\prod_{i \in I} a_i\right) = \prod_{i \in I} (1_i \star_i a_i) = \prod_{i \in I} a_i
    \]
    and
    \[
    	\left(\prod_{i \in I} a_i\right)\left(\prod_{i \in I} 1_i\right) = \prod_{i \in I} (a_i \star_i 1_i) = \prod_{i \in I} a_i.
    \]
    Thus, $\prod 1_i$ is the identity element of $G$.
    \item Let $\prod a_i$ be an element of $G$.
	Then
    \[
    	\left(\prod_{i \in I} a_i\right)\left(\prod_{i \in I} a_i^{-1}\right) = \prod_{i \in I} (a_i \star_i a_i^{-1}) = \prod_{i \in I} 1_i
    \]
    and
    \[
    	\left(\prod_{i \in I} a_i^{-1}\right)\left(\prod_{i \in I} a_i\right) = \prod_{i \in I} (a_i^{-1} \star_i a_i) = \prod_{i \in I} 1_i.
    \]
    Thus, $\prod a_i^{-1}$ is the inverse of $\prod a_i$.
	We conclude that $G$ is a group with inverse and identity as described above.
\end{solenum}

\begin{exercise}{5.1.16}
    State and prove the generalization of Proposition 2 to arbitrary direct products.
\end{exercise}

\begin{sol}
    Proposition 2 stated for direct arbitrary groups is as follows: if $G = \prod_{i \in I} G_i$ is a direct product of groups, then
    \begin{enumerate}
        \item Fix each $i$ with $j$-th coordinate where $j \neq i$ is the identity element of $G_j$ for $1 \leq j \leq n$.
	Then
        \[
        	\wh G_i = \set{(1, 1, \ldots, 1, g_i, 1, \ldots) \mid g_i \in G_i}
        \] 
        is a normal subgroup of $G$ that is isomorphic to $G_i$, or that
        \[
        	G/\wh G_i \cong \prod_{j \neq i} G_j.
        \]
        \item For fixed $i$, define $\pi_i : G \to G_i$ by $\pi_i((g_j)_{j \in I}) = g_i$.
	Then $\pi_i$ is a surjective homomorphism with kernel $\wh G_i$, so $G/\wh G_i \cong G_i$.
        \item If $x \in G_i$ and $y \in G_j$ for $i \neq j$, then $x$ and $y$ commute in $G$.
    \end{enumerate}
    The proof is similar to the proof of Proposition 2 in the text.
    \begin{enumerate}
        \item It is clear that $\wh G_i$ is a subgroup of $G$.
	Moreover, for any $\prod_{j \in I} g_j \in G$ and any $h = \prod_{j \neq i} 1_j h_i \in \wh G_i$, we have
        \[
        	\left(\prod_{j \in I} g_j\right)h\left(\prod_{j \in I} g_j\right)\inv = \prod_{j \neq i} 1_j (g_i h_i g_i\inv) \in \wh G_i
        \]
        since $g_i h_i g_i\inv \in G_i$.
	Thus, $\wh G_i$ is normal in $G$.
	Consider the map $\phi : G_i \to \wh G_i$ defined by $\phi(g_i) = \prod_{j \neq i} 1_j g_i$.
	It is clear that $\phi$ is a homomorphism.
	Moreover, if $\phi(g_i) = \prod_{j \neq i} 1_j g_i = \prod_{j \neq i} 1_j$, then $g_i$ is the identity element of $G_i$, so $\phi$ is injective.
	Finally, for any $\prod_{j \neq i} 1_j h_i \in \wh G_i$, we have $\phi(h_i) = \prod_{j \neq i} 1_j h_i$, so $\phi$ is surjective.
	Thus, $\wh G_i$ is isomorphic to $G_i$. 
        
        Consider the map
        \[
        	\psi : G \to \prod_{j \neq i} G_j \text{ defined by } \psi\left(\prod_{j \in I} g_j\right) = \prod_{j \neq i} g_j.
        \]
        It is clear that $\psi$ is a homomorphism.
	Moreover, the kernel of $\psi$ is $\wh G_i$, so by the First Isomorphism Theorem, we have
        \[
        	G/\wh G_i \cong \prod_{j \neq i} G_j.
        \]
        \item For fixed $i$, it is clear that $\pi_i$ is a homomorphism.
	Moreover, $\pi_i$ is surjective since for any $g_i \in G_i$, we have $\pi_i(\prod_{j \in I} 1_j g_i) = g_i$.
	Finally, the kernel of $\pi_i$ is $\wh G_i$, so by the First Isomorphism Theorem, we have
        \[
        	G/\wh G_i \cong G_i.
        \]
        \item Let $x = \prod_{j \neq i} 1_j x_i$ and $y = \prod_{j \neq k} 1_j y_k$ for some $i \neq k$.
	Then
        \[
        	xy = \prod_{j \neq i, k} 1_j x_i y_k = \prod_{j \neq i, k} 1_j y_k x_i = yx
        \]
        as desired.
    \end{enumerate}
\end{sol}

\begin{exercise}{5.1.17}
    Let $I$ be any nonempty index set and let $G_i$ be a group for each $i \in I$.
	The \emph{restricted direct product} or \emph{direct sum} of the groups $G_i$ is the set of elements of the direct product which are the identity in all but finitely many components, that is, the set of all elements $\prod a_i \in \prod_{i \in I} G_i$ such that $a_i = 1_i$ for all but a finite number of $i \in I$.
    \begin{subexercises}
        \item Prove that the restricted direct product is a subgroup of the direct product.
        \item Prove that the restricted direct product is normal in the direct product.
        \item Let $I = \zbb^+$ and let $p_i$ be the $i$th integer prime.
	    Show that if $G_i = \zbb/p_i\zbb$ for all $i \in \zbb^+$, then every element of the restricted direct product of the $G_i$'s has finite order but $\prod_{i \in \zbb^+} G_i$ has elements of infinite order.
	    Show that in this example the restricted direct product is the torsion subgroup of the direct product (cf. \exref{2.1.16}).
    \end{subexercises}
\end{exercise}

\begin{solenum}
    \item Let $G = \prod_{i \in I} G_i$ and let $H$ be the restricted direct product of the $G_i$'s.
	It is clear that $H$ is a subset of $G$.
	Moreover, if $\prod a_i$ and $\prod b_i$ are elements of $H$, then there are only finitely many $i$ such that $a_i \neq 1_i$ and only finitely many $i$ such that $b_i \neq 1_i$.
	Thus, there are only finitely many $i$ such that $a_i b_i \neq 1_i$, so $\prod a_i \prod b_i = \prod (a_i b_i)$ is an element of $H$.
	Finally, if $\prod a_i$ is an element of $H$, then there are only finitely many $i$ such that $a_i \neq 1_i$, so there are only finitely many $i$ such that $a_i^{-1} \neq 1_i$, so $\prod a_i^{-1}$ is an element of $H$.
	Thus, $H \leq G$.
    \item Let $G = \prod_{i \in I} G_i$ and let $H$ be the restricted direct product of the $G_i$'s.
	Let $\prod a_i$ be an element of $G$ and let $\prod b_i$ be an element of $H$.
	Then there are only finitely many $i$ such that $b_i \neq 1_i$, so there are only finitely many $i$ such that $a_i b_i a_i^{-1} \neq 1_i$, so $\prod a_i \prod b_i \prod a_i^{-1}$ is an element of $H$.
	Thus, $H \nsub G$.
    \item Let $G = \prod_{i \in \zbb^+} G_i$ where $G_i = \zbb/p_i\zbb$ for all $i$.
	Let $H$ be the restricted direct product of the $G_i$'s.
	If $\prod a_i$ is an element of $H$, then there are only finitely many $i$ such that $a_i \neq 1_i$, so there are only finitely many $i$ such that $a_i$ has order greater than 1, so $\prod a_i$ has finite order.
	On the other hand, if $\prod a_i$ is an element of $G$ such that $a_i$ is a generator of $G_i$ for all $i$, then $\prod a_i$ has infinite order since for any positive integer $n$, we have $(\prod a_i)^n = \prod a_i^n$, and there are infinitely many $i$ such that $a_i^n \neq 1_i$.
    
    To show $H = \tor(G)$, we note that $H \subseteq \tor(G)$ since every element of $H$ has finite order.
	Conversely, if $\prod a_i$ is an element of $\tor(G)$, then there is some positive integer $n$ such that $(\prod a_i)^n = \prod a_i^n = \prod 1_i$, so $a_i^n = 1_i$ for all $i$.
	Since $G_i$ is cyclic of prime order, this implies that $a_i$ is either the identity or a generator of $G_i$.
	If there are infinitely many $i$ such that $a_i$ is a generator of $G_i$, then $\prod a_i$ has infinite order, which is a contradiction.
	Thus, there are only finitely many $i$ such that $a_i$ is a generator of $G_i$, so there are only finitely many $i$ such that $a_i \neq 1_i$, so $\prod a_i$ is an element of $H$.
	Therefore, $\tor(G) \subseteq H$, and we conclude that $\tor(G) = H$.
\end{solenum}

\begin{exercise}{5.1.18}
    In each of (a) to (e) give an example of a group with the specified properties:
    \begin{subexercises}
        \item an infinite group in which every element has order 1 or 2
        \item an infinite group in which every element has finite order but for each positive integer $n$ there is an element of order $n$
        \item a group with an element of infinite order and an element of order 2
        \item a group $G$ such that every finite group is isomorphic to some subgroup of $G$
        \item a non-trivial group $G$ such that $G \cong G \times G$.
    \end{subexercises}
\end{exercise}

\begin{solenum}
    \item Consider the infinite direct product $\intmod[2] \times \intmod[2] \times \cdots$.
	Every element of this group is a tuple of 0's and 1's, and the order of each element is either 1 (if the tuple is the identity) or 2 (if the tuple has at least one 1).
    \item Let $G = \intmod[2] \times \intmod[3] \times \intmod[4] \times \cdots$.
	Every element of $G$ has finite order since it is a tuple of elements from finite cyclic groups.
	Moreover, for each positive integer $n$, the element that is the identity in all components except the $n$-th component, which is a generator of $\intmod[n]$, has order $n$.
    \item Consider the group $\intmod[2] \times \zbb$.
	The element $(1, 0)$ has order 2, and the element $(0, 1)$ has infinite order.
    \item By Cayley's Theorem, every finite group is isomorphic to a subgroup of some symmetric group $S_n$.
	We may then consider the group $G = S_1 \times S_2 \times S_3 \times \cdots$.
	For any finite group $H$, there is some $n$ such that $H$ is isomorphic to a subgroup of $S_n$, so $H$ is isomorphic to a subgroup of $G$.
    \item Let $G = \zbb \times \zbb \times \cdots$.
	Consider the map $\phi : G \to G \times G$ defined by
    \[
    	\phi((a_1, a_2, a_3, \ldots)) = ((a_1, a_3, a_5, \ldots), (a_2, a_4, a_6, \ldots)).
    \]
    It is clear that $\phi$ is a homomorphism.
	Moreover, if $\phi((a_1, a_2, a_3, \ldots)) = ((0, 0, 0, \ldots), (0, 0, 0, \ldots))$, then $a_i = 0$ for all $i$, so $\phi$ is injective.
	Finally, let $((b_1, b_2, b_3, \ldots), (c_1, c_2, c_3, \ldots))$ be an element of $G \times G$.
	Then 
    \[
    	\phi((b_1, c_1, b_2, c_2, b_3, c_3, \ldots)) = ((b_1, b_2, b_3, \ldots), (c_1, c_2, c_3, \ldots)),
    \]
    so $\phi$ is surjective.
	Thus, $\phi$ is an isomorphism and $G \cong G \times G$.
\end{solenum}